\uuid{fqJz}
\exo7id{576}
\titre{exo7 576}
\auteur{ridde}
\organisation{exo7}
\datecreate{1999-11-01}
\isIndication{false}
\isCorrection{false}
\chapitre{Suite}
\sousChapitre{Suites équivalentes, suites négligeables}
\cpge{mpsi-mp2i.nombres-reels-suites.c, pcsi.nombres-reels-suites.b, ptsi.nombres-reels-suites.b, mpsi-mp2i.nombres-reels-suites.e, pcsi.nombres-reels-suites.d, ptsi.nombres-reels-suites.d, mpsi-mp2i.nombres-reels-suites.d, pcsi.nombres-reels-suites.c, ptsi.nombres-reels-suites.c, bcpst1.suites-reelles}
\module{Analyse}
\niveau{L1}
\difficulte{}

\contenu{
\texte{
Soient $a_0$ et $b_0$ deux r\'eels fix\'es. On d\'efinit par r\'ecurrence les suites
$ (a_n)$ et $ (b_n)$ par $a_{n + 1} = \dfrac{2a_n + b_n}{3}$ et
 $b_{n + 1} = \dfrac{a_n + 2b_n}{3}$.
}
\begin{enumerate}
    \item \question{Montrer que ces deux suites sont adjacentes.}
    \item \question{En calculant $a_n + b_n$, montrer qu'elles convergent vers $\dfrac{a_0 + b_0}2$.}
\end{enumerate}
}
